% saida/questoes.tex

\begin{questao}{492177}
  \meta{tipo}{discursiva}
  \meta{banca}{Desconhecida}
  \meta{ano}{0}
  \meta{disciplina}{Física}
  \meta{topico}{Cinemática}
  \meta{conteudo}{Queda Livre}
  \meta{assunto}{Cálculo de tempo e posição}
  \meta{dificuldade}{medio}

  \enunciado{
    Em um experimento para se observar a queda dos corpos, dois sensores, A e B, foram abandonados de uma altura de 180 m sob ação exclusiva da aceleração da gravidade, cujo valor é 10 m/s$^2$. Ao tocar o solo, cada sensor envia um sinal que é detectado por um computador central que registra o tempo de queda. Dois segundos após abandonar o sensor A, larga-se o sensor B da mesma posição.

    \imagem{sensores.png}

    Após o sensor A ser abandonado, quanto tempo, em segundos, o computador leva para detectar o seu sinal?

    Quando o computador detecta o sinal do sensor A, em qual altura, em metros, estava o sensor B em relação ao solo?
  }

  \gabarito{}
\end{questao}

\begin{questao}{A-492177}
  \meta{tipo}{discursiva}
  \meta{banca}{Desconhecida}
  \meta{ano}{0}
  \meta{disciplina}{Física}
  \meta{topico}{Cinemática}
  \meta{conteudo}{Queda Livre}
  \meta{assunto}{Cálculo de tempo e posição}
  \meta{dificuldade}{facil}
  \meta{tags}{adaptada, NEE}

  \enunciado{
    \imagem{sensores.png}

    A imagem ilustra os sensores A e B sendo derrubados de uma altura.

    Dois sensores, A e B, foram deixados cair de uma altura de 180 metros. A aceleração da gravidade é de 10 m/s$^2$. O sensor B foi solto 2 segundos depois do A.

    \textbf{Quanto tempo o sensor A demorou para atingir o solo?}
  }

  \gabarito{}
\end{questao}

\begin{questao}{807361}
  \meta{tipo}{discursiva}
  \meta{banca}{Desconhecida}
  \meta{ano}{0}
  \meta{disciplina}{Física}
  \meta{topico}{Cinemática}
  \meta{conteudo}{Queda Livre}
  \meta{assunto}{Queda com intervalos iguais}
  \meta{dificuldade}{dificil}

  \enunciado{
    Uma torneira mal fechada pinga em intervalos de tempo iguais. A figura a seguir mostra a situação no instante em que uma das gotas está se soltando. Supondo que cada pingo abandone a torneira com velocidade nula e desprezando a resistência do ar, determine a razão A/B entre a distância A e B mostrada na figura.

    \imagem{torneira.png}
  }

  \gabarito{}
\end{questao}

\begin{questao}{A-807361}
  \meta{tipo}{discursiva}
  \meta{banca}{Desconhecida}
  \meta{ano}{0}
  \meta{disciplina}{Física}
  \meta{topico}{Cinemática}
  \meta{conteudo}{Queda Livre}
  \meta{assunto}{Queda com intervalos iguais}
  \meta{dificuldade}{medio}
  \meta{tags}{adaptada, NEE}

  \enunciado{
    \imagem{torneira.png}

    Na figura, há pingos de água caindo de uma torneira, e as distâncias entre eles são marcadas como A e B.

    Uma torneira está pingando, e cada pingo cai em intervalos de tempo iguais. Os pingos começam as suas quedas sempre a partir do repouso e sem a interferência da resistência do ar.

    \textbf{Qual é a razão entre as distâncias listadas na figura A e B?}
  }

  \gabarito{}
\end{questao}

\begin{questao}{186229}
  \meta{tipo}{discursiva}
  \meta{banca}{Desconhecida}
  \meta{ano}{0}
  \meta{disciplina}{Física}
  \meta{topico}{Cinemática}
  \meta{conteudo}{Queda Livre}
  \meta{assunto}{História da Física e Gravidade}
  \meta{dificuldade}{medio}

  \enunciado{
    \textbf{AGORA É COM @VOCÊ}

    O experimento científico de Galileu Galilei na Torre de Pisa é conhecido como um dos mais relevantes da história. Não há evidências concretas a respeito de sua realização, sendo a principal referência uma contestada citação na biografia de Galileu escrita por Vincenzo Viviani 12 anos após a morte do cientista. Entretanto, é um fato que se um experimento, que teria sido realizado em 1591, é descrito e comentado até hoje, isso mostra o quanto Galileu é importante para a história da Física e como as consequências desse hipotético experimento são significativas.

    Segundo relatos, Galileu Galilei soltou do alto da Torre de Pisa duas esferas com mesmo tamanho e massas diferentes. As esferas caíram com uma desprezível diferença no tempo de queda, contrariando a expectativa aristotélica, de que a esfera mais pesada cairia com maior velocidade. Com essas informações, que conclusão importante a respeito da queda dos corpos pode ser obtida?
  }

  \gabarito{}
\end{questao}

\begin{questao}{A-186229}
  \meta{tipo}{discursiva}
  \meta{banca}{Desconhecida}
  \meta{ano}{0}
  \meta{disciplina}{Física}
  \meta{topico}{Cinemática}
  \meta{conteudo}{Queda Livre}
  \meta{assunto}{História da Física e Gravidade}
  \meta{dificuldade}{facil}
  \meta{tags}{adaptada, NEE}

  \enunciado{
    O experimento científico de Galileu Galilei na Torre de Pisa é conhecido como um dos mais relevantes da história. Não há evidências concretas a respeito de sua realização, sendo a principal referência uma contestada citação na biografia de Galileu escrita por Vincenzo Viviani 12 anos após a morte do cientista. Entretanto, é um fato que se um experimento, que teria sido realizado em 1591, é descrito e comentado até hoje, isso mostra o quanto Galileu é importante para a história da Física e como as consequências desse hipotético experimento são significativas.

    O texto descreve um famoso experimento liderado por Galileu que mudou a história do pensamento científico.

    Galileu soltou do alto de uma torre duas esferas de massas diferentes, e ambas atingiram o solo num tempo de queda quase igual, o que foi contra a ideia de Aristóteles de que coisas mais pesadas cairiam mais rápido.

    \textbf{A partir das observações de Galileu, que conclusão podemos tirar a respeito da queda dos objetos?}
  }

  \gabarito{}
\end{questao}

\begin{questao}{669083}
  \meta{tipo}{discursiva}
  \meta{banca}{Desconhecida}
  \meta{ano}{0}
  \meta{disciplina}{Física}
  \meta{topico}{Cinemática}
  \meta{conteudo}{Queda Livre}
  \meta{assunto}{Cálculo de variação em altura}
  \meta{dificuldade}{medio}

  \enunciado{
    Uma bola de basquete cai em queda livre após ficar momentaneamente em repouso sobre o aro da tabela. Considerando que a altura do aro até o chão da quadra é de 3,2 m e que no local a aceleração da gravidade seja 10 m/s$^2$, responda ao que se pede.

    Calcule a velocidade, em m/s, com que a bola atinge o chão da quadra.

    Indique quanto tempo a bola leva, em segundos, para atingir o solo.
  }

  \gabarito{}
\end{questao}

\begin{questao}{A-669083}
  \meta{tipo}{discursiva}
  \meta{banca}{Desconhecida}
  \meta{ano}{0}
  \meta{disciplina}{Física}
  \meta{topico}{Cinemática}
  \meta{conteudo}{Queda Livre}
  \meta{assunto}{Cálculo de variação em altura}
  \meta{dificuldade}{facil}
  \meta{tags}{adaptada, NEE}

  \enunciado{
    Uma bola de basquete da altura do aro da tabela cai livremente para o chão perfeitamente livre de forças adversas. A altura do aro para o chão é de 3,2 m e a aceleração da gravidade atua a 10 m/s$^2$.

    \textbf{Tendo em vista os dados conhecidos do salto da bola de basquete, com que velocidade final ela atinge o solo?}
  }

  \gabarito{}
\end{questao}

\begin{questao}{470202}
  \meta{tipo}{discursiva}
  \meta{banca}{Desconhecida}
  \meta{ano}{0}
  \meta{disciplina}{Física}
  \meta{topico}{Cinemática}
  \meta{conteudo}{Queda Livre}
  \meta{assunto}{Intervalos de gravação}
  \meta{dificuldade}{medio}

  \enunciado{
    Um grupo de estudantes realizou um teste envolvendo queda livre no prédio da escola. Sob supervisão de seus professores e adotando procedimentos de segurança, como isolar e impedir a passagem de pessoas na parte inferior do prédio para evitar acidentes, os estudantes soltam, da base de uma janela, um objeto e tentam gravar a sua passagem pela janela do andar de baixo, onde se encontra uma câmera de vídeo que captura 120 quadros por segundo.

    \imagem{predio_janela.png}

    Em quantos quadros ocorre o registro da passagem do objeto? Considere g = 10 m/s$^2$.
  }

  \gabarito{}
\end{questao}

\begin{questao}{A-470202}
  \meta{tipo}{discursiva}
  \meta{banca}{Desconhecida}
  \meta{ano}{0}
  \meta{disciplina}{Física}
  \meta{topico}{Cinemática}
  \meta{conteudo}{Queda Livre}
  \meta{assunto}{Intervalos de gravação}
  \meta{dificuldade}{facil}
  \meta{tags}{adaptada, NEE}

  \enunciado{
    \imagem{predio_janela.png}

    A imagem revela a fachada de um prédio com estudantes realizando a captação através de câmeras nos andares de baixo e largando o objeto pelo andar de cima.

    Estudantes largaram um objeto de uma janela em queda livre com aceleração de 10 m/s$^2$. Na janela abaixo haverá uma câmera gravando que opera capturando 120 quadros em um único segundo.

    \textbf{Levando em consideração os elementos do sistema, em quantos quadros de gravação total ocorrerá o registro da passagem deste objeto?}
  }

  \gabarito{}
\end{questao}

\begin{questao}{277508}
  \meta{tipo}{discursiva}
  \meta{banca}{Desconhecida}
  \meta{ano}{0}
  \meta{disciplina}{Física}
  \meta{topico}{Cinemática}
  \meta{conteudo}{Lançamento Vertical}
  \meta{assunto}{Lançamento Vertical Ascendente}
  \meta{dificuldade}{medio}

  \enunciado{
    Um juiz de futebol lança verticalmente uma moeda para cima com velocidade de 2 m/s.

    Considerando g = 10 m/s$^2$, calcule

    o tempo de permanência da moeda no ar, em segundos, até voltar à mão do juiz.

    a altura máxima atingida pela moeda, em metros.
  }

  \gabarito{}
\end{questao}

\begin{questao}{A-277508}
  \meta{tipo}{discursiva}
  \meta{banca}{Desconhecida}
  \meta{ano}{0}
  \meta{disciplina}{Física}
  \meta{topico}{Cinemática}
  \meta{conteudo}{Lançamento Vertical}
  \meta{assunto}{Lançamento Vertical Ascendente}
  \meta{dificuldade}{facil}
  \meta{tags}{adaptada, NEE}

  \enunciado{
    Um juiz de futebol vai iniciar a partida efetuando o lançamento de uma moeda verticalmente para cima e esta tem velocidade inicial de 2 m/s. A gravidade age sobre o sistema a 10 m/s$^2$.

    \textbf{Qual valor, em metros, correspondente à altura de forma máxima alcançada pela moeda?}
  }

  \gabarito{}
\end{questao}

\begin{questao}{473383}
  \meta{tipo}{objetiva}
  \meta{banca}{UPF}
  \meta{ano}{0}
  \meta{disciplina}{Física}
  \meta{topico}{Cinemática}
  \meta{conteudo}{Lançamento Vertical}
  \meta{assunto}{Análise Gráfica (v x t)}
  \meta{dificuldade}{medio}

  \enunciado{
    (UPF) Durante um experimento, um estudante joga uma bola de tênis verticalmente para cima e observa sua trajetória até retornar à sua mão.

    Desprezando a resistência do ar, qual gráfico melhor representa a velocidade (v) da bola em relação ao tempo de movimento (t) durante todo o percurso?
  }

  \begin{alternativas}
    \alt{A}{\imagem{grafico_t_1.png}}
    \alt{B}{\imagem{grafico_t_2.png}}
    \alt{C}{\imagem{grafico_t_3.png}}
    \alt{D}{\imagem{grafico_t_4.png}}
    \alt{E}{\imagem{grafico_t_5.png}}
  \end{alternativas}

  \gabarito{D}
\end{questao}

\begin{questao}{A-473383}
  \meta{tipo}{objetiva}
  \meta{banca}{UPF}
  \meta{ano}{0}
  \meta{disciplina}{Física}
  \meta{topico}{Cinemática}
  \meta{conteudo}{Lançamento Vertical}
  \meta{assunto}{Análise Gráfica (v x t)}
  \meta{dificuldade}{facil}
  \meta{tags}{adaptada, NEE}

  \enunciado{
    Durante um projeto de experimento, lançou-se uma bola de tênis verticalmente em direção ao céu e ela posteriormente caiu retornando às mãos do seu lançador sem perder energia.

    \textbf{Tendo em mente a relação de um gráfico em plano cartesiano ligando as propriedades de velocidade em função do tempo, como deverá ser a sua representação?}
  }

  \begin{alternativas}
    \alt{A}{\imagem{grafico_t_1.png}}
    \alt{B}{\imagem{grafico_t_4.png}}
    \alt{C}{\imagem{grafico_t_5.png}}
  \end{alternativas}

  \gabarito{B}
\end{questao}

\begin{questao}{174024}
  \meta{tipo}{discursiva}
  \meta{banca}{Desconhecida}
  \meta{ano}{0}
  \meta{disciplina}{Física}
  \meta{topico}{Cinemática}
  \meta{conteudo}{Lançamento Vertical}
  \meta{assunto}{Tempo de Voo}
  \meta{dificuldade}{medio}

  \enunciado{
    \textbf{AGORA É COM @VOCÊ}

    Para tentar descobrir a altura de uma árvore, dois amigos fazem o seguinte teste: um deles jogou uma pequena bola para cima até a altura aproximada do topo da árvore, enquanto o segundo mediu o tempo total para a bola subir até o topo e retornar à mão de quem a lançou. Sendo 1,8 s o tempo total que a bola ficou no ar, desprezando-se a resistência do ar e considerando a aceleração da gravidade igual a 10 m/s$^2$, responda ao que se pede.

    Qual a altura da árvore?

    Com qual velocidade a bola foi lançada?
  }

  \gabarito{}
\end{questao}

\begin{questao}{A-174024}
  \meta{tipo}{discursiva}
  \meta{banca}{Desconhecida}
  \meta{ano}{0}
  \meta{disciplina}{Física}
  \meta{topico}{Cinemática}
  \meta{conteudo}{Lançamento Vertical}
  \meta{assunto}{Tempo de Voo}
  \meta{dificuldade}{facil}
  \meta{tags}{adaptada, NEE}

  \enunciado{
    Dois amigos jogam uma bola verticalmente tentando chegar perfeitamente ao topo de uma árvore. O tempo total da bola (subida e a volta até sua mão) foi pontuado em 1,8 segundos contando com que a gravidade atua a 10 m/s$^2$.

    \textbf{Identificando essas características do movimento, qual tamanho da altura final foi atingida pela bola de modo paralelo a árvore?}
  }

  \gabarito{}
\end{questao}

\begin{questao}{811666}
  \meta{tipo}{discursiva}
  \meta{banca}{Desconhecida}
  \meta{ano}{0}
  \meta{disciplina}{Física}
  \meta{topico}{Cinemática}
  \meta{conteudo}{Queda Livre}
  \meta{assunto}{Encontro de Corpos}
  \meta{dificuldade}{dificil}

  \enunciado{
    Do alto de uma plataforma de mergulho, duas pedras são lançadas para baixo. Uma delas foi solta a partir do repouso e atinge a piscina depois de 2 s. A segunda, lançada com velocidade inicial não nula, atinge a piscina depois de 1,8 s. Desprezando-se a resistência do ar e considerando g = 10 m/s$^2$, qual o valor da velocidade de lançamento, em m/s, da segunda pedra?
  }

  \gabarito{}
\end{questao}

\begin{questao}{A-811666}
  \meta{tipo}{discursiva}
  \meta{banca}{Desconhecida}
  \meta{ano}{0}
  \meta{disciplina}{Física}
  \meta{topico}{Cinemática}
  \meta{conteudo}{Queda Livre}
  \meta{assunto}{Encontro de Corpos}
  \meta{dificuldade}{medio}
  \meta{tags}{adaptada, NEE}

  \enunciado{
    Existem duas pedras que são lançadas para a água, do alto da plataforma. A inicial foi perfeitamente solta do repouso e atingiu a piscina em exatos 2 segundos, mas a segunda teve lançamento direto durando apenas 1,8 segundos e a aceleração aplicada de base é a da Terra que é equivalente a 10 m/s$^2$.

    \textbf{Tendo em vista a descida da segunda pedra, em qual valor de medida em m/s encontra-se a velocidade da impulsão deste lançamento?}
  }

  \gabarito{}
\end{questao}

\begin{questao}{565377}
  \meta{tipo}{discursiva}
  \meta{banca}{Desconhecida}
  \meta{ano}{0}
  \meta{disciplina}{Física}
  \meta{topico}{Cinemática}
  \meta{conteudo}{Lançamento Vertical}
  \meta{assunto}{Cálculos em outros planetas}
  \meta{dificuldade}{medio}

  \enunciado{
    Uma empresa de tecnologia aeroespacial está projetando um dispositivo que irá recolher amostras da atmosfera de Marte. O dispositivo lança um sensor verticalmente para cima com certa velocidade inicial. No ponto mais alto atingido pelo sensor, um pequeno paraquedas se abre e algumas medições são realizadas durante a queda. Primeiramente, o dispositivo foi testado na superfície da Terra com a mesma velocidade inicial projetada para o lançamento a ser realizado em Marte.

    Qual deve ser a razão entre a altura máxima esperada para ser alcançada em Marte e a altura obtida na Terra? Para fins práticos, despreze a resistência do ar nos dois planetas e considere a gravidade na Terra como 10 m/s$^2$ e a de Marte igual a 4 m/s$^2$.
  }

  \gabarito{}
\end{questao}

\begin{questao}{A-565377}
  \meta{tipo}{discursiva}
  \meta{banca}{Desconhecida}
  \meta{ano}{0}
  \meta{disciplina}{Física}
  \meta{topico}{Cinemática}
  \meta{conteudo}{Lançamento Vertical}
  \meta{assunto}{Cálculos em outros planetas}
  \meta{dificuldade}{facil}
  \meta{tags}{adaptada, NEE}

  \enunciado{
    Um dispositivo lançou verticalmente um dispositivo de forma rápida na Terra para testar as suas subidas, e efetuará o exato mesmo lançamento, em um futuro próximo de uma viagem direcionada dentro do planeta Marte. A resistência sendo ignorada e a gravidade dos planetas sendo respectivamente a da Terra sendo equivalente a 10 m/s$^2$ e a de Marte tendo no total prático 4 m/s$^2$.

    \textbf{Para prever de forma conclusiva a viagem, qual o cálculo da razão da altura a se formar em Marte para a testada inicialmente na Terra?}
  }

  \gabarito{}
\end{questao}

\begin{questao}{646462}
  \meta{tipo}{objetiva}
  \meta{banca}{UEA}
  \meta{ano}{2022}
  \meta{disciplina}{Física}
  \meta{topico}{Cinemática}
  \meta{conteudo}{Queda Livre}
  \meta{assunto}{Ausência de Resistência do Ar}
  \meta{dificuldade}{facil}

  \enunciado{
    \textbf{ATIVIDADES PROPOSTAS}

    A Nasa Space Power Facility é um laboratório que possui a maior câmara de vácuo do mundo. No interior dessa câmara foi conduzido um experimento utilizando uma bola de boliche e uma pena. Ambas foram erguidas a uma mesma altura, e a câmara foi evacuada. Após esse procedimento, em ausência de ar, a bola de boliche e a pena foram abandonadas simultaneamente.

    Ao final desse experimento, observou-se que
  }

  \begin{alternativas}
    \alt{A}{a bola atingiu o solo primeiro por ter massa maior que a da pena, mesmo com a ausência de ar.}
    \alt{B}{a pena permaneceu parada onde estava e a bola atingiu o solo.}
    \alt{C}{a pena e a bola adquiriram uma aceleração, durante a queda, menor do que a da gravidade local.}
    \alt{D}{a pena, por ser menor que a bola, atingiu o solo primeiro.}
    \alt{E}{a bola e a pena atingiram o solo simultaneamente.}
  \end{alternativas}

  \gabarito{E}
\end{questao}

\begin{questao}{A-646462}
  \meta{tipo}{objetiva}
  \meta{banca}{UEA}
  \meta{ano}{2022}
  \meta{disciplina}{Física}
  \meta{topico}{Cinemática}
  \meta{conteudo}{Queda Livre}
  \meta{assunto}{Ausência de Resistência do Ar}
  \meta{dificuldade}{facil}
  \meta{tags}{adaptada, NEE}

  \enunciado{
    A instituição Nasa Space Power Facility contém a gigantesca forma da maior câmara laboratorial de um enorme vácuo contido e lacrado. Se inserirmos para fins de experimento uma bola pesada de boliche e logo ao lado colocarmos uma leve pena e abandonarmos essas ao longo desse ambiente sem nenhum ar presente, ambas caem sobre o mesmo efeito daquela determinada ausência ali.

    \textbf{Caso sejam largadas a uma mesma altura ao mesmo tempo no vácuo o que deve ser observado de modo conclusivo sobre o alcance ao chão dos itens citados?}
  }

  \begin{alternativas}
    \alt{A}{A bola e a leve pena chegam perfeitamente conjuntas atuando ambas com sucesso em bater o fundo simultaneamente.}
    \alt{B}{A bola vai ser puxada pra atingir o local no fundo em tempo rápido dado da força de sua própria capacidade e a sua massa pesada.}
    \alt{C}{A pena ficará planando e flutuante pela condição e sem chance de bater ao piso que era determinado alvo principal.}
  \end{alternativas}

  \gabarito{A}
\end{questao}

\begin{questao}{114539}
  \meta{tipo}{objetiva}
  \meta{banca}{PUC-RJ}
  \meta{ano}{2023}
  \meta{disciplina}{Física}
  \meta{topico}{Cinemática}
  \meta{conteudo}{Lançamento Vertical}
  \meta{assunto}{Velocidade Inicial}
  \meta{dificuldade}{medio}

  \enunciado{
    Uma bola de massa 2,0 kg é lançada verticalmente para cima, a partir do solo. Após 0,5 s, sua velocidade é a metade daquela de lançamento. Com qual velocidade, em m/s, a bola é lançada?

    \textbf{Dado:} g = 10 m/s$^2$.
  }

  \begin{alternativas}
    \alt{A}{2,0}
    \alt{B}{5,0}
    \alt{C}{10}
    \alt{D}{20}
    \alt{E}{50}
  \end{alternativas}

  \gabarito{C}
\end{questao}

\begin{questao}{A-114539}
  \meta{tipo}{objetiva}
  \meta{banca}{PUC-RJ}
  \meta{ano}{2023}
  \meta{disciplina}{Física}
  \meta{topico}{Cinemática}
  \meta{conteudo}{Lançamento Vertical}
  \meta{assunto}{Velocidade Inicial}
  \meta{dificuldade}{facil}
  \meta{tags}{adaptada, NEE}

  \enunciado{
    Lança-se no céu e sempre vindo de baixo, uma bola da massa com 2,0 kg e então após se passar no decorrer do tempo 0,5 segundos ela possui em mãos exatamente a porção reduzida à sua metade natural provinda dos momentos originais desta largada. A gravidade aplica os mesmos 10 m/s$^2$ normais.

    \textbf{Com a informação deste impacto que alterou na sua fase de diminuição, por qual marco estava avaliada em inicial este arremesso total e respectiva velocidade do corpo?}
  }

  \begin{alternativas}
    \alt{A}{5,0 m/s}
    \alt{B}{20 m/s}
    \alt{C}{10 m/s}
  \end{alternativas}

  \gabarito{C}
\end{questao}

\begin{questao}{224383}
  \meta{tipo}{objetiva}
  \meta{banca}{FGV}
  \meta{ano}{0}
  \meta{disciplina}{Física}
  \meta{topico}{Cinemática}
  \meta{conteudo}{Lançamento Vertical}
  \meta{assunto}{Ponto Mais Alto}
  \meta{dificuldade}{medio}

  \enunciado{
    Uma pedra foi arremessada verticalmente para cima, a partir do solo, e permaneceu no ar por 4 s até regressar ao solo. Desprezando a resistência do ar e adotando g = 10 m/s$^2$, a altura máxima atingida por essa pedra foi
  }

  \begin{alternativas}
    \alt{A}{5 m.}
    \alt{B}{10 m.}
    \alt{C}{15 m.}
    \alt{D}{20 m.}
    \alt{E}{25 m.}
  \end{alternativas}

  \gabarito{D}
\end{questao}

\begin{questao}{A-224383}
  \meta{tipo}{objetiva}
  \meta{banca}{FGV}
  \meta{ano}{0}
  \meta{disciplina}{Física}
  \meta{topico}{Cinemática}
  \meta{conteudo}{Lançamento Vertical}
  \meta{assunto}{Ponto Mais Alto}
  \meta{dificuldade}{facil}
  \meta{tags}{adaptada, NEE}

  \enunciado{
    O objeto sólido (a pedra) permaneceu ao voo pelo tempo contado por cravados 4 segundos após um arremesso direto, logo em seguida sem atritos extras na sua jornada tocou de novo sobre o gramado sob influência final aplicada aos usuais fatores com aceleração 10 m/s$^2$.

    \textbf{Tirando da relação dessa metade focada sobre um só eixo, até qual quantidade dos metros ela atinge e demonstra o próprio momento liminar e altíssimo em sua máxima de voo?}
  }

  \begin{alternativas}
    \alt{A}{20 m.}
    \alt{B}{10 m.}
    \alt{C}{25 m.}
  \end{alternativas}

  \gabarito{A}
\end{questao}

\begin{questao}{794086}
  \meta{tipo}{objetiva}
  \meta{banca}{UTFPR}
  \meta{ano}{0}
  \meta{disciplina}{Física}
  \meta{topico}{Cinemática}
  \meta{conteudo}{Queda Livre}
  \meta{assunto}{Análise Gráfica}
  \meta{dificuldade}{medio}

  \enunciado{
    Assinale o gráfico que melhor representa a velocidade em função do tempo durante uma queda livre desprezando a resistência do ar.
  }

  \begin{alternativas}
    \alt{A}{\imagem{grafico_v_1.png}}
    \alt{B}{\imagem{grafico_v_2.png}}
    \alt{C}{\imagem{grafico_v_3.png}}
    \alt{D}{\imagem{grafico_v_4.png}}
    \alt{E}{\imagem{grafico_v_5.png}}
  \end{alternativas}

  \gabarito{D}
\end{questao}

\begin{questao}{A-794086}
  \meta{tipo}{objetiva}
  \meta{banca}{UTFPR}
  \meta{ano}{0}
  \meta{disciplina}{Física}
  \meta{topico}{Cinemática}
  \meta{conteudo}{Queda Livre}
  \meta{assunto}{Análise Gráfica}
  \meta{dificuldade}{facil}
  \meta{tags}{adaptada, NEE}

  \enunciado{
    Para verificar formas e também traços variados e precisos contidos no padrão para objetos arremessados num trajeto de queda do repouso sem perder ar.

    \textbf{Conferindo o fato para todo movimento e por causa dos limites lineares constantes associados no impacto da velocidade, o que seria o correto de acordo para descrever esta função através e diretamente ao meio do papel em gráfico?}
  }

  \begin{alternativas}
    \alt{A}{\imagem{grafico_v_1.png}}
    \alt{B}{\imagem{grafico_v_4.png}}
    \alt{C}{\imagem{grafico_v_5.png}}
  \end{alternativas}

  \gabarito{B}
\end{questao}

\begin{questao}{399485}
  \meta{tipo}{objetiva}
  \meta{banca}{UNISINOS}
  \meta{ano}{0}
  \meta{disciplina}{Física}
  \meta{topico}{Cinemática}
  \meta{conteudo}{Queda Livre}
  \meta{assunto}{Cinemática Mista}
  \meta{dificuldade}{dificil}

  \enunciado{
    Uma pequena pedra é abandonada do repouso do alto de uma ponte, 45 m acima do nível da água de um rio que passa por baixo dela. A pedra cai dentro de uma pequena embarcação que se move com velocidade constante e que estava a 12 m de distância do ponto de impacto, no instante em que a pedra foi solta. O módulo da velocidade com que a embarcação se move no rio é

    \textbf{Dado:} g = 10 m/s$^2$.
  }

  \begin{alternativas}
    \alt{A}{1 m/s.}
    \alt{B}{2 m/s.}
    \alt{C}{3 m/s.}
    \alt{D}{4 m/s.}
    \alt{E}{5 m/s.}
  \end{alternativas}

  \gabarito{D}
\end{questao}

\begin{questao}{A-399485}
  \meta{tipo}{objetiva}
  \meta{banca}{UNISINOS}
  \meta{ano}{0}
  \meta{disciplina}{Física}
  \meta{topico}{Cinemática}
  \meta{conteudo}{Queda Livre}
  \meta{assunto}{Cinemática Mista}
  \meta{dificuldade}{medio}
  \meta{tags}{adaptada, NEE}

  \enunciado{
    Ao cair abandonado ao nível repousado numa ponte, esse pequeno e simples pedaço de mineral teve de descer cerca de quarenta e cinco (45 m) rumando reto e profundo da maré que estava passando ali. E nisso existia ali embaixo ainda uma condução de barco percorrendo a trechos distantes, e incrivelmente ela apossava na mesma marca dos 12 e sendo com aceleração livre na base nos mesmos padrões em 10 m/s$^2$.

    \textbf{Com todo esse grande choque com essa rocha pequena dentro dele, de quanto representava propriamente a velocidade total contida para deslocamento que essa tal embarcação e modesta canoa fazia navegando no mar?}
  }

  \begin{alternativas}
    \alt{A}{1 m/s.}
    \alt{B}{4 m/s.}
    \alt{C}{5 m/s.}
  \end{alternativas}

  \gabarito{B}
\end{questao}

\begin{questao}{888861}
  \meta{tipo}{objetiva}
  \meta{banca}{UERJ}
  \meta{ano}{0}
  \meta{disciplina}{Física}
  \meta{topico}{Cinemática}
  \meta{conteudo}{Queda Livre}
  \meta{assunto}{Proporcionalidade com g}
  \meta{dificuldade}{medio}

  \enunciado{
    O Sistema Solar é formado por planetas que apresentam diferentes acelerações da gravidade. Admita que um corpo é solto em queda livre na Terra a uma altura \textbf{h} e atinge a superfície do planeta com velocidade de 5 m/s. Admita ainda um planeta P, também do Sistema Solar, em que o mesmo corpo é solto, à mesma altura \textbf{h}, e atinge velocidade final de 8 m/s.

    Sabe-se que o quadrado da velocidade com a qual um corpo em queda livre atinge a superfície é diretamente proporcional à aceleração da gravidade do planeta. Considere os valores aproximados apresentados na tabela.

    \begin{tabular}{|l|c|}    \hline
    \textbf{Planeta} & \textbf{Aceleração da gravidade (m/s$^2$)} \\
    \hline
    Júpiter & 25 \\
    Marte & 4 \\
    Netuno & 11 \\
    Terra & 10 \\
    Vênus & 9 \\
    \hline
    \end{tabular}

    Com base nessas informações, o planeta que apresenta a aceleração da gravidade mais próxima à do planeta P é
  }

  \begin{alternativas}
    \alt{A}{Júpiter.}
    \alt{B}{Marte.}
    \alt{C}{Netuno.}
    \alt{D}{Vênus.}
  \end{alternativas}

  \gabarito{A}
\end{questao}

\begin{questao}{A-888861}
  \meta{tipo}{objetiva}
  \meta{banca}{UERJ}
  \meta{ano}{0}
  \meta{disciplina}{Física}
  \meta{topico}{Cinemática}
  \meta{conteudo}{Queda Livre}
  \meta{assunto}{Proporcionalidade com g}
  \meta{dificuldade}{facil}
  \meta{tags}{adaptada, NEE}

  \enunciado{
    \begin{tabular}{|l|c|}    \hline
    \textbf{Planeta} & \textbf{Aceleração da gravidade (m/s$^2$)} \\
    \hline
    Júpiter & 25 \\
    Marte & 4 \\
    Netuno & 11 \\
    Terra & 10 \\
    Vênus & 9 \\
    \hline
    \end{tabular}

    A tabela apresenta planetas do Sistema Solar e quanto vale a sua respectiva e real da atração gravitacional (agindo direto para o centro).

    Uma observação feita compara ao mundo misterioso e em um planeta ocultado “P” o teste que foi verificado ser o seu trajeto batendo sobre do planeta com fortes resultados para a frente equivalentes da batida e resultando de 8 e na Terra este teste resultaria só de apenas e pequenos 5 de valor.

    \textbf{Para prever do que foi citado nos elementos e com aprovação em análises feitas, entre estes que constam nos citados qual deverá vir do mais e mais próximo e exato e também verossímil a respeito aos da gravidade do tão estimado planeta do P?}
  }

  \begin{alternativas}
    \alt{A}{Marte.}
    \alt{B}{Júpiter.}
    \alt{C}{Vênus.}
  \end{alternativas}

  \gabarito{B}
\end{questao}

\begin{questao}{708737}
  \meta{tipo}{objetiva}
  \meta{banca}{FMC}
  \meta{ano}{0}
  \meta{disciplina}{Física}
  \meta{topico}{Cinemática}
  \meta{conteudo}{Queda Livre}
  \meta{assunto}{Percursos iguais}
  \meta{dificuldade}{dificil}

  \enunciado{
    Uma pedra é abandonada na beira de um abismo e, desprezando-se a resistência do ar, cai sob a ação apenas da força gravitacional. Dois segundos após, uma segunda pedra é abandonada da mesma posição e, nesse instante, a distância entre as duas pedras é \textbf{d}. Dois segundos após a segunda pedra ser abandonada, a distância entre elas será
  }

  \begin{alternativas}
    \alt{A}{d.}
    \alt{B}{2d.}
    \alt{C}{3d.}
    \alt{D}{4d.}
    \alt{E}{5d.}
  \end{alternativas}

  \gabarito{C}
\end{questao}

\begin{questao}{A-708737}
  \meta{tipo}{objetiva}
  \meta{banca}{FMC}
  \meta{ano}{0}
  \meta{disciplina}{Física}
  \meta{topico}{Cinemática}
  \meta{conteudo}{Queda Livre}
  \meta{assunto}{Percursos iguais}
  \meta{dificuldade}{medio}
  \meta{tags}{adaptada, NEE}

  \enunciado{
    Duas pedras minúsculas são postas perto a frente à margem livre de abismos as atirado ao limbo separadamente mas juntas do mesmo localizador e tendo exatamente nos 2 segundos sendo um diferencial marcando e nomeando o espaço dessa distância provida nas separadas de uma para outras pelas distâncias marcadas sendo chamada só letra "d".

    \textbf{Passado esse momento após as suas saídas com os relógios pontuados em mais extras 2 segundos como resultará o que será deixado para suas grandezas espaçadas da distância da uma até às outras no meio do final das contas?}
  }

  \begin{alternativas}
    \alt{A}{2d.}
    \alt{B}{3d.}
    \alt{C}{5d.}
  \end{alternativas}

  \gabarito{B}
\end{questao}

\begin{questao}{378638}
  \meta{tipo}{objetiva}
  \meta{banca}{UECE}
  \meta{ano}{0}
  \meta{disciplina}{Física}
  \meta{topico}{Cinemática}
  \meta{conteudo}{Queda Livre}
  \meta{assunto}{Propriedades MRUV}
  \meta{dificuldade}{medio}

  \enunciado{
    Um objeto em queda livre percorre uma distância X durante o primeiro segundo de queda. Desprezando as forças resistivas e considerando a aceleração da gravidade constante ao longo de todo o percurso, a distância percorrida por esse objeto durante o quarto segundo da queda corresponde a
  }

  \begin{alternativas}
    \alt{A}{3X.}
    \alt{B}{5X.}
    \alt{C}{4X.}
    \alt{D}{7X.}
  \end{alternativas}

  \gabarito{D}
\end{questao}

\begin{questao}{A-378638}
  \meta{tipo}{objetiva}
  \meta{banca}{UECE}
  \meta{ano}{0}
  \meta{disciplina}{Física}
  \meta{topico}{Cinemática}
  \meta{conteudo}{Queda Livre}
  \meta{assunto}{Propriedades MRUV}
  \meta{dificuldade}{facil}
  \meta{tags}{adaptada, NEE}

  \enunciado{
    Encontra-se a grandeza misteriosa e chamada com nome de X em um espaço físico, marcando um distanciamento nos testes aéreos que medem o quanto foi voado já dentro no seu tempo com exato começo e a entrada para seu primeiro de modo a se basear todo o seu seguimento natural nas métricas ideais e perfeitamente no vácuo de forma mantida da aceleração.

    \textbf{Qual medida representaria o fato real no espaço para os momentos no momento do andamento no seu próprio decorrer de quarto (4) de segundo para este corpo percorrido do decorrer?}
  }

  \begin{alternativas}
    \alt{A}{3X.}
    \alt{B}{4X.}
    \alt{C}{7X.}
  \end{alternativas}

  \gabarito{C}
\end{questao}

\begin{questao}{620573}
  \meta{tipo}{objetiva}
  \meta{banca}{EEAR}
  \meta{ano}{0}
  \meta{disciplina}{Física}
  \meta{topico}{Cinemática}
  \meta{conteudo}{Lançamento Vertical}
  \meta{assunto}{Equações Horárias}
  \meta{dificuldade}{dificil}

  \enunciado{
    Um professor cronometra o tempo “$t_S$” que um objeto (considerado um ponto material) lançado a partir do solo, verticalmente para cima e com uma velocidade inicial, leva para realizar um deslocamento $\Delta x_S$ até atingir a altura máxima. Em seguida, o professor mede, em relação à altura máxima, o deslocamento de descida $\Delta x_D$ ocorrido em um intervalo de tempo igual a 1/4 de “$t_S$” cronometrado inicialmente.

    \textbf{Dado:} Considere o módulo da aceleração da gravidade constante e que, durante todo o movimento do objeto, não há nenhum tipo de atrito.

    A razão $\frac{\Delta x_S}{\Delta x_D}$ é igual a
  }

  \begin{alternativas}
    \alt{A}{2.}
    \alt{B}{4.}
    \alt{C}{8.}
    \alt{D}{16.}
  \end{alternativas}

  \gabarito{D}
\end{questao}

\begin{questao}{A-620573}
  \meta{tipo}{objetiva}
  \meta{banca}{EEAR}
  \meta{ano}{0}
  \meta{disciplina}{Física}
  \meta{topico}{Cinemática}
  \meta{conteudo}{Lançamento Vertical}
  \meta{assunto}{Equações Horárias}
  \meta{dificuldade}{medio}
  \meta{tags}{adaptada, NEE}

  \enunciado{
    Existem cronômetros no poder do professor e o mesmo usou as táticas medindo tempo do "tS", esse no ar voo que o deslocou até parar na "altura máxima" ou também apelidado sendo no escrito matemático seu deslocar com "$\Delta x_S$". Logo assim no retorno de descida livre, há agora ele tirando do limite um medidor da metade das parciais das quedas onde tem nos tempos em cerca sendo equivalentes à sua forma na fração e com formato e fração de um quarto (1/4) desse seu começo já percorrido.

    \textbf{Tendo uma ausência do ar totalizada e considerando suas metades que foram dadas e do formato dos valores de proporção descrita a divisão inteira desta conta vira ao real ou a sua razão sendo?}
  }

  \begin{alternativas}
    \alt{A}{4.}
    \alt{B}{8.}
    \alt{C}{16.}
  \end{alternativas}

  \gabarito{C}
\end{questao}

\begin{questao}{951094}
  \meta{tipo}{objetiva}
  \meta{banca}{UPE}
  \meta{ano}{0}
  \meta{disciplina}{Física}
  \meta{topico}{Cinemática}
  \meta{conteudo}{Lançamento Vertical e Queda Livre}
  \meta{assunto}{Encontro de corpos}
  \meta{dificuldade}{dificil}

  \enunciado{
    Do alto de um edifício de altura 20,0 m, deixou-se cair, desde o repouso, uma pequena esfera de aço; ao mesmo tempo, no chão desse edifício, na mesma linha imaginária, lançou-se para cima outra esfera de aço com velocidade $v_0$. Qual é aproximadamente a velocidade $v_0$, em m/s, da segunda esfera para que elas se encontrem na meia altura do edifício?

    \textbf{Dados:} g = 10 m/s$^2$ e $\sqrt{2}$ = 1,4.
  }

  \begin{alternativas}
    \alt{A}{4}
    \alt{B}{10}
    \alt{C}{14}
    \alt{D}{20}
    \alt{E}{25}
  \end{alternativas}

  \gabarito{C}
\end{questao}

\begin{questao}{A-951094}
  \meta{tipo}{objetiva}
  \meta{banca}{UPE}
  \meta{ano}{0}
  \meta{disciplina}{Física}
  \meta{topico}{Cinemática}
  \meta{conteudo}{Lançamento Vertical e Queda Livre}
  \meta{assunto}{Encontro de corpos}
  \meta{dificuldade}{medio}
  \meta{tags}{adaptada, NEE}

  \enunciado{
    Pronto ao chão fica o edifício grande totalizado sendo tendo de cume de altura medido tendo exatos seus 20m completos deixando com repouso nulo soltar os do teto e lançando nos andares da calçada logo de encontro sendo a colisão entre objetos como em formato redondo feitos por ligas vindos da mistura de metais que a velocidade inicial deve providenciar exatamente no cume na altura da área da média inteira destas bases cruzadas. As considerações básicas vêm dos conhecidos e base na fórmula do dez na atração com dez m/s$^2$ assim que a raiz exata base das duas metades chegue e aproxime num dos $1,4$. 

    \textbf{Por fim, determinando de certo do valor na taxa aproximada da quantia exigida, como fará a medida m/s que dará exato choque para a sua metade ideal do prédio para o andamento?}
  }

  \begin{alternativas}
    \alt{A}{4 m/s.}
    \alt{B}{14 m/s.}
    \alt{C}{25 m/s.}
  \end{alternativas}

  \gabarito{B}
\end{questao}
